\section{Introduction}
Diffusion models have become one of the main paradigms for high-quality image generation, with successful architectures ranging from U-Net denoisers to latent diffusion models and diffusion transformers~\cite{ronneberger2015unet,ho2020ddpm,rombach2022ldm,peebles2023dit}. However, training such models remains computationally expensive: not only must the model learn the denoising dynamics, but it must also be able to extract useful internal visual representations (discriminative features) in its hidden states~\cite{li2023dreamteacher,xiang2023ddae,chen2024deconstructing,mukhopadhyay2023diffusion}

Recent work on REPA~\cite{yu2025repa} proposes to accelerate diffusion training by explicitly aligning intermediate hidden states of the denoising model with clean image representations extracted by a frozen pretrained visual encoder. This suggests a simple but powerful idea: instead of forcing the diffusion model to discover good internal representations only through the denoising loss, we can use an external representation model as an auxiliary teacher.

In this project, we study representation alignment as a training principle for efficient diffusion modeling under constrained computational settings. We focus on reduced datasets, low-data regimes, and fine-grained visual domains. At first glance, one might expect REPA to be especially useful in such settings, since the pretrained teacher can inject visual knowledge that is unavailable from the small training set alone. However, this advantage is not obvious. Recent evidence from iREPA~\cite{singh2026irepa} suggests that the effectiveness of REPA is determined not merely by the global quality of the teacher encoder, but by the spatial structure of its features. Therefore, in narrow domains, a teacher encoder may provide globally meaningful but locally mismatched representations, which can make the alignment objective less helpful or even harmful.

Our goal is to characterize when REPA improves convergence speed and sample quality, how its effectiveness depends on the pretrained feature extractor and dataset scale, and under which conditions the alignment objective can become a limitation rather than a reliable source of inductive bias. We also will try to verify the usefulness of REPA loss for text-diffusion models, as recently shown by~\cite{junior2026accelerating}.
